% Modernised reading — Grothendieck fonds, folder 115, batch 1 (leaves 1-14).
%
% This is an INTERPRETATION, not a transcription. It reads the leaves in
% today's notation and today's names, states hypotheses the manuscript leaves
% implicit, and drops the critical apparatus so the mathematics can be read
% continuously. Everything the brackets and underlines used to carry — an
% uncertain reading, a notation he used differently, a passage that is simply
% missing — is said in a footnote instead.
%
% It is held to being mathematically correct as it stands, not merely faithful
% to the leaf. Where the manuscript is loose or the reading is doubtful, the
% statement here is the one that is true, and the footnote says what the leaf
% actually has. Where the two disagree in substance, the transcription is the
% record: anyone citing Grothendieck should cite batch-01.fr.tex.
%
% Pass: Fable 5 (claude-fable-5), 2026-08-08 — first pass, unchecked against
% the leaves by a human.
\documentclass[11pt,a4paper]{article}
% Le préambule commun à toutes les transcriptions du fonds Grothendieck.
%
% Il tient en une idée : **l'incertitude se compose, elle ne se résout pas.**
% Une transcription qui lisse un mot douteux a détruit la seule chose pour
% laquelle on la faisait — un lecteur doit pouvoir savoir, à chaque endroit,
% ce qui était sur le feuillet et ce qui a été ajouté. Les six macros qui
% suivent sont donc l'appareil critique, et non de la décoration.
%
% Les mêmes commandes sont reconnues par `scripts/render.mjs`, qui produit la
% vue de lecture du site. Une macro ajoutée ici doit l'être aussi là-bas,
% faute de quoi le rendu échouera bruyamment — ce qui est voulu : un rendu
% silencieusement faux est pire qu'un rendu absent.

\ProvidesPackage{grothendieck}[2026/08/08 Transcriptions du fonds Grothendieck]

\usepackage{amsmath,amssymb,amsthm}
\usepackage{mathrsfs}
\usepackage{stmaryrd}
% Les diagrammes commutatifs sont partout dans ce fonds : c'est la langue
% même de la géométrie algébrique de Grothendieck, et une transcription qui
% les rendrait en prose ne transcrirait rien.
\usepackage{tikz-cd}
% Français par défaut — c'est la langue des feuillets — mais l'anglais est
% chargé aussi : la traduction et le résumé sont des documents anglais, et la
% typographie française (espaces avant les deux-points) y serait une faute.
% Ces documents basculent par \selectlanguage{english} en tête de corps.
\usepackage[english,french]{babel}
\usepackage{fontspec}
\usepackage{geometry}
\geometry{a4paper,margin=25mm,marginparwidth=28mm}
\usepackage{marginnote}
\usepackage{xcolor}
% ulem, not soul. `soul`'s \st and \ul are built on a character-by-character
% scan that XeLaTeX's font machinery breaks: the document fails with an
% undefined control sequence rather than degrading. ulem does the same two jobs
% and is Unicode-engine safe. `normalem` stops it hijacking \emph, which the
% transcriptions use for Grothendieck's own emphasis.
\usepackage[normalem]{ulem}
\usepackage{hyperref}
\hypersetup{colorlinks=true,linkcolor=black,urlcolor=black,pdfborder={0 0 0}}

% --- Métadonnées du lot -----------------------------------------------------
% Reprises telles quelles par le script de rendu, qui les affiche en tête de
% la vue de lecture. Elles fixent ce que le fichier prétend couvrir : sans
% elles, une transcription flotte sans point d'attache dans un fonds de seize
% mille feuillets.

\newcommand{\folder}[1]{\gdef\grothFolder{#1}}
\newcommand{\batch}[1]{\gdef\grothBatch{#1}}
\newcommand{\leaves}[2]{\gdef\grothFirst{#1}\gdef\grothLast{#2}}
\newcommand{\foldertitle}[1]{\gdef\grothFoldertitle{#1}}
\gdef\grothFolder{?}\gdef\grothBatch{?}\gdef\grothFirst{?}\gdef\grothLast{?}\gdef\grothFoldertitle{}

% --- Le résumé --------------------------------------------------------------
%
% Il ouvre la lecture modernisée, sous le titre « Résumé », et il est composé
% en petit. Ce n'est pas un ornement : c'est le seul passage du document qui
% ne soit pas des mathématiques, et un lecteur doit pouvoir le reconnaître
% avant de l'avoir lu — puis le sauter s'il connaît déjà le sujet, ou s'y
% arrêter s'il ne le connaît pas. Le filet qui le suit marque la reprise du
% corps.

\newenvironment{resume}
  {\small\setlength{\parskip}{0.45em}}
  {\par\medskip\hrule\medskip}

% --- Mots-clés ---------------------------------------------------------------
%
% En anglais, en clôture du résumé de la lecture modernisée : le vocabulaire
% moderne sous lequel chercher ce que les feuillets construisent. Le manifeste
% du site (`npm run manifest`) les extrait pour étiqueter la cote entière —
% cette ligne est la source unique des tags, il n'y a pas de fichier à côté.
\newcommand{\keywords}[1]{\par\smallskip\noindent
  {\footnotesize\textsc{Keywords} --- \textit{#1}}\par}

% --- L'appareil critique ----------------------------------------------------

% Le numéro de feuillet, en marge. C'est l'ancre qui permet de régler un
% désaccord sur une formule en regardant *un* feuillet plutôt qu'une chemise
% de six cents. Le site s'en sert aussi pour tourner les pages du fac-similé
% quand on fait défiler la transcription.
\newcommand{\leaf}[1]{%
  \marginnote{\footnotesize\color{gray!70}\textsf{#1}}%
  \label{leaf:#1}\ignorespaces}

% Illisible : on ne devine pas. Un mot inventé qui se lit comme les autres est
% le pire résultat possible.
\newcommand{\ill}{\textcolor{red!60!black}{[\,\ldots\,]}}

% Lecture douteuse : le mot est donné, et signalé comme incertain.
\newcommand{\uncertain}[1]{\uline{#1}}

% Ajout de l'éditeur — une lettre manquante, un mot sous-entendu.
\newcommand{\add}[1]{\textcolor{blue!50!black}{[#1]}}

% Biffé par Grothendieck lui-même. Ce qu'il a rayé fait partie du document :
% une rature dit souvent où la pensée a changé de direction.
\newcommand{\struck}[1]{\sout{#1}}

% Note du transcripteur, sur l'état matériel du feuillet.
\newcommand{\note}[1]{\par\smallskip{\small\color{gray!60!black}\textit{[#1]}}\par\smallskip}

% Note marginale de Grothendieck, distincte du corps du texte.
\newcommand{\marginal}[1]{\par\smallskip\noindent\hspace{1em}%
  {\small\color{gray!60!black}#1}\par\smallskip}

% --- Titre ------------------------------------------------------------------

\AtBeginDocument{%
  \begin{center}
    {\large\bfseries Fonds Alexandre Grothendieck --- cote n\textsuperscript{o}~\grothFolder}\\[2pt]
    {\itshape\grothFoldertitle}\\[4pt]
    {\small feuillets \grothFirst--\grothLast\ (lot \grothBatch)}\\[2pt]
    {\footnotesize\color{gray!60!black}Transcription établie d'après le fac-similé numérisé
     par l'Université de Montpellier.\\
     Les lectures incertaines sont soulignées, les illisibles notées [\,\ldots\,],
     les ajouts entre crochets.}
  \end{center}
  \vspace{1em}}

\endinput


\folder{115}
\batch{1}
\leaves{1}{14}
\foldertitle{« Correspondances » fonctorielles. Dualité des topos — modernised reading}

\begin{document}
\selectlanguage{english}

\section*{Isbell duality, as folder 115 works it out}

Throughout, $A$ and $B$ are small categories. Write
$\widehat{A} = [A^{\mathrm{op}}, \mathbf{Set}]$ for presheaves and
$A^{\vee} = [A, \mathbf{Set}]$ for copresheaves.\footnote{Grothendieck's own
notation, stable across the folder. He writes $\mathrm{Ens}$ for
$\mathbf{Set}$; we do not.} Both are presheaf categories — $A^{\vee} =
\widehat{A^{\mathrm{op}}}$ — hence both are Grothendieck topoi, which is what
the folder's title refers to. Write $y$ for the Yoneda embedding.

Two decorations recur. A subscript shriek, $\underline{\mathrm{Hom}}_{!}$,
denotes the full subcategory of functors preserving colimits; a superscript
one, $\underline{\mathrm{Hom}}^{!}$, those preserving limits; a doubled shriek
the two-variable version, separately in each argument.\footnote{Leaf 3, head
of the leaf. The manuscript says \emph{limites inductives} and \emph{limites
projectives}, both marked as uncertain readings in the transcription; the
equivalences that follow hold only under this assignment, which is what fixes
it.} In current usage: cocontinuous and continuous.

\subsection*{Kernels between presheaf topoi}

The first movement is a calculus of correspondences, on the analogy of an
integral kernel $k(x,y)$ carrying functions on one space to functions on
another.

Let $p_{A}, p_{B}$ be the projections off $A \times B$, and let the relative
hom mean the pushforward of the internal hom along the projection that
forgets the remaining variable.\footnote{The manuscript's subscripts are the
opposite of the modern reflex: his $p_{B}$ is the projection that
\emph{deletes} $B$, so $p_{B} : A \times B \to A$. Read the other way the
identities on leaf 3 contradict one another. This is the one point in the
batch where the reading turns on a convention rather than on a letterform.}
For $X \in (A \times B)^{\wedge}$ the two transforms are
\[
{}^{s}X(L) = \underline{\mathrm{Hom}}_{B}(p_{B}^{*}L,\ X), \qquad
{}^{d}X(M) = \underline{\mathrm{Hom}}_{A}(p_{A}^{*}M,\ X),
\qquad L \in \widehat{A},\ M \in \widehat{B},
\]
and they are related by the single pairing
\[
X(L,M) = \mathrm{Hom}(L \boxtimes M,\ X)
\;\cong\; \mathrm{Hom}_{\widehat{A}}\bigl(L,\ {}^{d}X(M)\bigr)
\;\cong\; \mathrm{Hom}_{\widehat{B}}\bigl(M,\ {}^{s}X(L)\bigr),
\qquad L \boxtimes M = p_{B}^{*}L \times p_{A}^{*}M .
\]
The external product $\boxtimes : \widehat{A} \times \widehat{B} \to
(A \times B)^{\wedge}$ is fully faithful.

Note the variance, because it is the whole point of what follows. Since
$(A\times B)^{\wedge} = [A^{\mathrm{op}} \times B^{\mathrm{op}},
\mathbf{Set}]$, the kernel $X$ is contravariant in \emph{both} arguments, so
${}^{s}X$ and ${}^{d}X$ are contravariant functors
\[
{}^{s}X : \widehat{A}^{\,\mathrm{op}} \longrightarrow \widehat{B},
\qquad
{}^{d}X : \widehat{B}^{\,\mathrm{op}} \longrightarrow \widehat{A},
\]
and the displayed isomorphisms make them \emph{mutually right adjoint} — an
adjunction between $\widehat{A}^{\,\mathrm{op}}$ and $\widehat{B}$. That is a
dual adjunction, not an ordinary one.\footnote{So $X$ is not a profunctor
$A \nrightarrow B$: a profunctor is covariant in one variable. It is a
two-variable presheaf, and the construction it generates is a Galois
connection of the sort $V \mapsto V^{*}$ is for vector spaces — which is
exactly why the folder arrives where it does.}

Leaf 5 arranges eight functor categories in a wheel, every arrow of which is
an equivalence, among them
\[
(A\times B)^{\wedge} \;\simeq\;
\underline{\mathrm{Hom}}(A^{\mathrm{op}}, \widehat{B}) \;\simeq\;
\underline{\mathrm{Hom}}^{!}(\widehat{A}^{\,\mathrm{op}}, \widehat{B}) \;\simeq\;
\underline{\mathrm{Hom}}_{!}(A^{\vee}, \widehat{B}).
\]
Each is a way of currying the same two-variable functor. The first is plain
currying; the second and third extend along Yoneda, and the shriek records
what the extension must preserve for it to be an equivalence — a functor out
of a free (co)completion is determined by its restriction to the generators
exactly when it preserves the (co)limits that completion adds.\footnote{The
eight arise from choosing which variable is taken outside, and at which of
the three levels — $A^{\mathrm{op}}$, $\widehat{A}^{\,\mathrm{op}}$,
$A^{\vee}$ — it is presented, the extremes closing on the two-variable forms
$\underline{\mathrm{Hom}}^{!!}$ and $\underline{\mathrm{Hom}}_{!!}$. The
transcription marks the wheel's two left-hand nodes as uncertain, a tear
running through that side of the leaf; they are recoverable from the symmetry,
and are given here accordingly.} The same statement is then carried to finite
families $(A_{i})_{i \in I}$ with $I$ split into a covariant and a
contravariant part.

\subsection*{The duality}

Everything above specialises. Take $B = A^{\mathrm{op}}$ and for the kernel
take the hom itself, $H_{A} = \mathrm{Hom}_{A} \in (A \times
A^{\mathrm{op}})^{\wedge}$ — the manuscript's own remark, in the left margin
of leaf 3, that this is the canonical element of that category. The dual
adjunction of the previous section becomes an adjunction between $\widehat{A}$
and $(A^{\vee})^{\mathrm{op}}$, and this is Isbell duality.

Explicitly, its two functors are
\[
\mathcal{O}(\Phi)(a) = \widehat{A}\bigl(\Phi,\ y a\bigr), \qquad
\mathrm{Spec}(\Psi)(a) = A^{\vee}\bigl(\Psi,\ \mathrm{Hom}_{A}(a,-)\bigr),
\]
covariant and contravariant in $a$ respectively, so that $\mathcal{O}\Phi \in
A^{\vee}$ and $\mathrm{Spec}\,\Psi \in \widehat{A}$; and
\[
H_{A}(\Phi, \Psi)
\;\cong\; \mathrm{Hom}_{\widehat{A}}\bigl(\Phi,\ \mathrm{Spec}\,\Psi\bigr)
\;\cong\; \mathrm{Hom}_{A^{\vee}}\bigl(\Psi,\ \mathcal{O}\Phi\bigr),
\]
which is $\mathcal{O} \dashv \mathrm{Spec}$ as functors $\widehat{A}
\rightleftarrows (A^{\vee})^{\mathrm{op}}$.\footnote{The manuscript writes
$\varphi_{A}$ for $\mathcal{O}$ and $\Psi_{A}$ for $\mathrm{Spec}$; the names
$\mathcal{O}$ and $\mathrm{Spec}$ are later, and carry the geometric reading
of the adjunction as functions-and-points. Nothing in the folder cites
Isbell's work of the mid-1960s, and the construction is arrived at from the
kernel calculus above rather than from it.}

Each functor is determined up to unique isomorphism by two properties, and
the manuscript states both: it restricts along $y$ to the covariant hom, and
it carries colimits in its source to limits in its target.\footnote{The
manuscript says $\Psi_{A}$ ``commutes with projective limits''. Read in
$(A^{\vee})^{\mathrm{op}}$ that is the statement above: $\mathrm{Spec}$ is a
right adjoint, so it preserves limits of $(A^{\vee})^{\mathrm{op}}$, which are
colimits of $A^{\vee}$.} On representables the pairing is the hom:
$H_{A}(y a, \mathrm{Hom}_{A}(b,-)) \cong \mathrm{Hom}_{A}(a,b)$, since
$\mathcal{O}(y a) = \mathrm{Hom}_{A}(a,-)$ and Yoneda does the rest.

The unit and counit
\[
\eta_{\Phi} : \Phi \longrightarrow \mathrm{Spec}\,\mathcal{O}\,\Phi,
\qquad
\varepsilon_{\Psi} : \Psi \longrightarrow \mathcal{O}\,\mathrm{Spec}\,\Psi
\]
are therefore isomorphisms on representables. Whether they are isomorphisms
anywhere else is the question the folder poses and leaves open:

\begin{quote}
Are the adjunction morphisms isomorphisms only for representables? Is it a
weak equivalence — probably not.
\end{quote}

This is the reflexivity question of linear algebra one level up. For vector
spaces over a field, $V \to V^{**}$ is an isomorphism exactly on the
finite-dimensional ones; here the objects at which $\eta$ is invertible are
the \emph{reflexive}, or Isbell-self-dual, objects, and they are what the next
section constructs.\footnote{His parenthesis ``sans doute pas'' is retained
above; the transcription marks the words before it as an uncertain reading.}

Dually there is a contraction, cocontinuous in each variable,
\[
H^{A} : \widehat{A} \times A^{\vee} \longrightarrow \mathbf{Set},
\qquad
\Phi \star_{A} \Psi \;=\; \int^{a \in A} \Phi(a) \times \Psi(a),
\]
the coend, agreeing with the hom on representables by co-Yoneda, and it gives
\[
\widehat{A} \;\simeq\; \underline{\mathrm{Hom}}_{!}(A^{\vee}, \mathbf{Set}),
\qquad
A^{\vee} \;\simeq\; \underline{\mathrm{Hom}}_{!}(\widehat{A}, \mathbf{Set}),
\]
each topos being the cocontinuous functionals on the other.\footnote{Both are
instances of the free cocompletion property: $\widehat{A}$ is the free
cocompletion of $A$, so cocontinuous functors out of it correspond to functors
out of $A$. The decorations on the second differ from the first in the
manuscript and are flagged uncertain in the transcription; they are given here
as matching, which is what the symmetry requires.}

\subsection*{The envelope}

Leaves 10 and 11 build what is now called the Isbell envelope. Let $A
\hookrightarrow B$ be full and $x$ an object of $B$. The two homs $H_{x} =
\mathrm{Hom}_{B}(-, x) \in \widehat{A}$ and $H^{x} = \mathrm{Hom}_{B}(x, -)
\in A^{\vee}$ compose, by composition in $B$:
\[
H_{x}(a) \times H^{x}(b) \xrightarrow{\ \alpha_{x}\ } \mathrm{Hom}_{A}(a,b),
\]
which lands in $\mathrm{Hom}_{A}$ rather than $\mathrm{Hom}_{B}$ precisely
because $A$ is full. Abstracting from $B$, let $\widetilde{A}$ be the category
of triples $(H_{*}, H^{*}, \alpha)$ with $H_{*} \in \widehat{A}$, $H^{*} \in
A^{\vee}$ and $\alpha$ natural. A morphism is a pair
\[
\varphi_{*} : H_{*} \to H'_{*} \ \text{in } \widehat{A},
\qquad
\varphi^{*} : H'^{*} \to H^{*} \ \text{in } A^{\vee},
\]
covariant in one component and contravariant in the other, compatible with the
pairings:
\[
\alpha_{a,b}\bigl(u,\ \varphi^{*}(u')\bigr)
= \alpha'_{a,b}\bigl(\varphi_{*}(u),\ u'\bigr).
\]
Every full embedding $A \hookrightarrow B$ then induces $B \to \widetilde{A}$;
the construction is functorial in $A$; and it is canonically self-dual,
$\widetilde{A}^{\,\mathrm{op}} \simeq \widetilde{A^{\mathrm{op}}}$, the
duality simply exchanging the two components.\footnote{The functoriality
square closing leaf 11 has several arrow labels marked conjectural in the
transcription; the shape is unambiguous where the labels are not.} Both
$\widehat{A}$ and $(A^{\vee})^{\mathrm{op}}$ embed in it, by
$\Phi \mapsto (\Phi, \mathcal{O}\Phi, \mathrm{ev})$ and dually — which is the
sense in which $\widetilde{A}$ is where the two sides of the duality meet.

Leaf 13 places the Cauchy completion at the centre. The identity
\[
\bigl(\mathrm{Kar}(A^{\mathrm{op}})\bigr)^{\mathrm{op}} \;\simeq\; \mathrm{Kar}(A)
\]
holds because splitting idempotents commutes with passage to the opposite, and
for $\mathbf{Set}$-enriched categories $\mathrm{Kar}$ — the Karoubi envelope —
\emph{is} the Cauchy completion. Retracts of representables are reflexive, so
$\mathrm{Kar}(A)$ sits inside the self-dual part of
$\widetilde{A}$.\footnote{The manuscript writes $A = \widehat{A} \cap
A^{\vee\circ}$, an intersection of subcategories of different categories,
which cannot be read literally; and equality with $\mathrm{Kar}(A)$ does not
hold in general — the reflexive objects can be strictly more numerous. The
containment is what is safe to assert, and is what the rest of the leaf uses.}

The leaf closes with a caution about gluing, worth keeping because it is a
warning and not a step: given $N = M_{1} \cap M_{2}$ through full embeddings,
one may form $\Sigma'$ from $M'_{1}$ and $M'_{2}$ over it, but even when
$M'_{i} \to M_{i}$ are equivalences the images need not be full in $\Sigma'$,
and $N' \to N$ need not be an equivalence.\footnote{Several words of this
passage are illegible in the manuscript. The negations are legible, and they
are what carries the warning.}

\subsection*{Coefficients}

The last leaf lets the duality take values in a category $M$ rather than in
$\mathbf{Set}$. For $M$ cocomplete,
\[
\underline{\mathrm{Hom}}(A, M) \;\simeq\; \underline{\mathrm{Hom}}_{!}(\widehat{A}, M),
\]
the universal property of $\widehat{A}$ as the free cocompletion of
$A$;\footnote{The cocompleteness of $M$ is not stated on the leaf. Without it
the right-hand side is not defined, so it is supplied here.} and, for
$M$ moreover complete,
\[
\mathcal{F}(X, M)^{\mathrm{op}} \;\simeq\; \check{\mathcal{F}}(X, M^{\mathrm{op}}),
\]
exchanging sheaves and cosheaves under passage to the opposite. The leaf
breaks off there.\footnote{Its decorations are the least legible in the batch,
and the transcription marks most of the final formulas as uncertain. They are
given here in the only form consistent with the rest of the folder.}

\end{document}
